\section{The normalized cubic theta input and its radial profiles}
\label{sec:theta}

To prove the reflected moment, we first realize the exterior coefficient
\(\gamma_2(\mathfrak c)/(q_{\mathfrak c}^{1/2}q_{\mathfrak n})\)
in the family \(B_m\).  The cubic theta representation used here has
one-dimensional local Whittaker spaces.  Its selected central datum,
together with normalized induction, gives the factor \(q_{\mathfrak p}^{-1}\)
for a triple valuation step.  Compatible rational, torus, root, and compact
sections fix the unit factor at valuation \(e\) as
\(\chi_{\mathfrak p}(u)^{2e}\).  We state this normalized input first and
then identify its exterior coefficient.  The verification of the
Kazhdan--Patterson specialization is in \Cref{app:whittaker}.

Throughout this section \(F\) is the fixed cyclotomic field from the
preceding section, and \(\Adele_F\) is its adele ring.  In particular,
\(\mu_{12}\subset F\).  We use the standard trace character
\(\psi:F\backslash\Adele_F\longrightarrow\C^\times\), and choose the
self-dual additive Haar measures, for which
\(F\backslash\Adele_F\) has volume one.  Write
\[
 n(x)=\begin{pmatrix}1&x\\0&1\end{pmatrix},\qquad
 \bar n(x)=\begin{pmatrix}1&0\\x&1\end{pmatrix},\qquad
 a(y)=\begin{pmatrix}y&0\\0&1\end{pmatrix},\qquad
 w=\begin{pmatrix}0&1\\1&0\end{pmatrix}.
\]
The same notation will denote the specified lifts below.  A Whittaker
functional with character \(\psi_v\) means a linear functional \(\lambda\)
satisfying
\(\lambda(\pi(n(x))\varphi)=\psi_v(x)\lambda(\varphi)\).

\subsection{The automorphic input and its exterior coefficients}

We use the determinant-modified cover of Kazhdan and Patterson with
parameter \(c=2\), with the ordinary torus section and the canonical root
splittings.  The embedding \(\iota:\mu_3(F)\hookrightarrow\C^\times\)
is oriented as in the exterior cocycle formula below.  If \(a_v(y)\)
denotes the chosen torus lift, define the effective scalar
\(\mathcal C_v(x,y)\) by
\[
 \pi(a_v(x))\pi(a_v(y))
   =\mathcal C_v(x,y)\pi(a_v(xy))
\]
in a genuine representation \(\pi\).  Thus this notation includes the
chosen embedding of the central \(\mu_3\).

\begin{proposition}[Cubic theta input]\label{prop:theta-input}
For the fixed field \(F\), there is a threefold central cover
\[
 1\longrightarrow\mu_3
 \longrightarrow\widetilde{\PGL}_2(\Adele_F)
 \longrightarrow\PGL_2(\Adele_F)\longrightarrow1
\]
which splits over \(\PGL_2(F)\), and an irreducible genuine automorphic
representation
\(\Theta=\bigotimes'_v\Theta_v\) on this cover, with the following
properties.

\begin{enumerate}[label=\textup{(\roman*)}]
\item \emph{Global representation and Fourier expansion.}
  The cover is the quotient of the \(c=2\) cover
  \(\widetilde{\GL}_2^{(2)}\) by the ordinary split scalar subgroup.
  The inducing character of \(\Theta_v\), on the center of the covering
  torus before this quotient, is
  \begin{equation}\label{eq:theta-central-datum}
  \omega_v\!\left(\zeta\,s_v(\operatorname{diag}(x^3,y^3))s_v(zI)\right)
   =\iota(\zeta)\left|x^3/y^3\right|_v^{1/6}.
  \end{equation}
  Here \(s_v\) is the ordinary torus section.  The normalized inducing
  norm exponents are \(1/6,-1/6\).  The local exceptional representation is the irreducible
  quotient for these exponents and the irreducible subrepresentation for
  their Weyl conjugates.
  Every \(\Theta_v\) has a one-dimensional Whittaker space.
  The global functional
  \[
   \lambda(\varphi)=
    \int_{F\backslash\Adele_F}\theta_\varphi(n(x))
          \overline{\psi(x)}\,\dd x
  \]
  on the automorphic realization \(\varphi\mapsto\theta_\varphi\) is
  nonzero.  The representation is spherical at all but finitely many
  finite places.
  There is a preliminary pure tensor \(\varphi^\circ\) on which
  \(\lambda\) is nonzero.  Its finitely many nondistinguished finite
  places lie in the final fixed \(S\), along with the archimedean places,
  the places over \(6\), and the places where \(\psi_v\) has nonzero conductor.

  For a pure tensor
  \(\varphi=\varphi_S\otimes\bigotimes_{v\notin S}\varphi_v\),
  with the distinguished spherical vector at almost every \(v\notin S\),
  the global Whittaker function factors as
  \begin{equation}\label{eq:theta-normalized-factorization}
  \mathcal W_\varphi(g)
   :=\lambda(\pi(g)\varphi)
   =\mathcal W_{S,\varphi_S}(g_S)
       \prod_{v\notin S}W_{v,\varphi_v}(g_v).
  \end{equation}
  Each exterior spherical functional is normalized at one, so only
  finitely many factors in this product differ from one for adelic \(g\).
  The functional at \(S\) factors at its complex places; at its finite
  places it may be retained as a finite sum of products of local Whittaker
  functionals.  For a smooth automorphic vector,
  \begin{equation}\label{eq:theta-fourier-expansion}
   \theta_\varphi(g)=\theta_{\varphi,N}(g)
    +\sum_{\nu\in F^\times}\mathcal W_\varphi(a(\nu)g),\qquad
   \theta_{\varphi,N}(g)
    =\int_{F\backslash\Adele_F}\theta_\varphi(n(x)g)\,\dd x .
  \end{equation}
  This expansion converges absolutely and locally uniformly on compact
  sets of \(g\).

\item \emph{Exact lifts and exterior local values.}
  On the quotient the chosen lifts satisfy
  \begin{equation}\label{eq:theta-lift-identities}
  \begin{gathered}
   n(x)a(y)=a(y)n(x/y),\qquad
   wa(y)w^{-1}=a(y^{-1}),\\
   wn(H)=n(H^{-1})a(-H^{-2})\bar n(H^{-1}).
  \end{gathered}
  \end{equation}
  In the last identity \(H\ne0\).  At an exterior prime, when \(H^{-1}\)
  is integral the last factor is compact, and a compact diagonal may absorb
  the minus sign.

  For every exterior prime \(\mathfrak p\), choose a uniformizer
  \(\varpi\), put \(q_{\mathfrak p}=\Norm\mathfrak p\), and use the
  compact splitting on \(\PGL_2(\mathcal O_{\mathfrak p})\).
  The orientation of \(\iota\), equivalently that of the Hilbert-symbol
  convention, can be fixed so that the scalar by which the torus
  cocycle acts in a genuine representation is
  \begin{equation}\label{eq:theta-effective-cocycle}
   \mathcal C_{\mathfrak p}(\varpi,u)=\chi_{\mathfrak p}(u)^{-2}
       \qquad(u\in\mathcal O_{\mathfrak p}^{\times}).
  \end{equation}
  The compact splitting agrees with the chosen sections on unit
  diagonals, Weyl elements, and integral upper and lower root elements.
  Define
  \[
   g_{2,\mathfrak p}
    =q_{\mathfrak p}^{-1/2}
       \sum_{x\bmod\mathfrak p}
       \chi_{\mathfrak p}(x)^2\psi_{\mathfrak p}(x/\varpi).
  \]
  Then \(|g_{2,\mathfrak p}|=1\).  If \(W_{\mathfrak p}\) is the
  spherical Whittaker function normalized by \(W_{\mathfrak p}(1)=1\),
  then, for \(u\in\mathcal O_{\mathfrak p}^{\times}\),
  \begin{equation}\label{eq:theta-local-values}
  W_{\mathfrak p}(a(\varpi^e u))
   =\chi_{\mathfrak p}(u)^{2e}
   \begin{cases}
    0,&e<0,\\
    q_{\mathfrak p}^{-k},&e=3k,\quad k\ge0,\\
    q_{\mathfrak p}^{-k-1/2}g_{2,\mathfrak p},
        &e=3k+1,\quad k\ge0,\\
    0,&e=3k+2,\quad k\ge0.
   \end{cases}
  \end{equation}
  In particular a shift by three nonnegative valuations multiplies the
  value by \(q_{\mathfrak p}^{-1}\).

\item \emph{Complex local behavior.}
  At a complex place, \(\Theta_v\) is the ordinary spherical
  \(\PGL_2(\C)\) principal series with normalized characters
  \((|\cdot|_v^{1/6},|\cdot|_v^{-1/6})\), where \(|y|_v\) is the square
  of the ordinary modulus.  Its spherical Whittaker function is a
  nonzero constant times
  \begin{equation}\label{eq:theta-complex-bessel}
   |y|_v^{1/2}K_{1/3}\bigl(c_v|y|_v^{1/2}\bigr),\qquad c_v>0.
  \end{equation}
  The constant term of the automorphic theta function belongs to the
  Weyl-conjugate torus induction.  In particular its dependence on a
  complex diagonal \(a(y)\) is radial, also when the vector on its right
  has been acted on by a fixed compact element.
\end{enumerate}
\end{proposition}

In \eqref{eq:theta-normalized-factorization}, the condition
\(W_{\mathfrak p}(1)=1\) fixes every exterior spherical scalar; the
remaining normalization belongs to the \(S\)-functional.
\Cref{app:whittaker} derives this product from the nonzero global
functional and local uniqueness, including the correction to the original
global formula.  The noncompact norm
integrations below require the separate absolute bounds in
\Cref{lem:sheet-profile}.

\begin{corollary}[The exterior theta coefficient]
\label{cor:theta-exterior-coefficient}
At every exterior prime use an allowed prime generator as the uniformizer
in \Cref{prop:theta-input}.  Let \(A\in R\setminus\{0\}\) be an allowed
generator of an integral exterior ideal \(\mathfrak A\).  Then
\begin{equation}\label{eq:theta-exterior-coefficient}
 \prod_{\mathfrak p\notin S}W_{\mathfrak p}(a(A))
 =
 \begin{cases}
 \displaystyle
 \frac{\gamma_2(\mathfrak c)}
      {q_{\mathfrak c}^{1/2}q_{\mathfrak n}},
   &\mathfrak A=\mathfrak c\mathfrak n^3,\quad
     \mathfrak c\ \text{squarefree},\\[2mm]
 0,&\text{otherwise}.
 \end{cases}
\end{equation}
In the first case \(\mathfrak c,\mathfrak n\) are uniquely determined;
they need not be coprime.
\end{corollary}

\begin{proof}
The local support in \eqref{eq:theta-local-values} permits precisely the
valuations \(3k\) and \(3k+1\).  They give the unique decomposition
\(\mathfrak A=\mathfrak c\mathfrak n^3\), with \(\mathfrak c\) squarefree,
and magnitudes \(q_{\mathfrak p}^{-k}\) and
\(q_{\mathfrak p}^{-k-1/2}\), respectively.  Use allowed product
generators \(A=cn^3\); any other allowed \(A\) differs by an element of
\(U_S^6\), which has no effect on the local values.  At a prime
\(\mathfrak p\mid\mathfrak c\), the local unit factor is
\(\chi_{\mathfrak p}(c/p)^2\).  The unit from \(n^3\) has exponent
divisible by six in the symbol, including when \(\mathfrak p\mid\mathfrak n\),
and disappears.  The normalized prime sum \(g_{2,\mathfrak p}\) equals
\(\gamma_2(\mathfrak p)\): the local additive sign from
\Cref{eq:normalized-gauss} has no effect because
\(\chi_{\mathfrak p}(-1)=1\).  The product of the unit factors and these
prime sums is \(\gamma_2(\mathfrak c)\) by \Cref{eq:gauss-crt}.
Multiplying the magnitudes gives \eqref{eq:theta-exterior-coefficient}.
\end{proof}

Thus theta supplies the exterior factor in \(B_m\).  The ideal character
will enter through an \(S\)-average, and finite upper-unipotent translations
will insert the moving sextic twist.  The chosen \(S\)-Whittaker function
first supplies a radial profile; an exact radial average will replace it
by \(P_G\).  We prepare those \(S\)-components next.

The arithmetic power classes also describe the torus factor at \(S\).
Put \(\mathcal C_S=\prod_{v\in S}\mathcal C_v\).  At an exterior prime,
\Cref{eq:theta-effective-cocycle} and the tame Hilbert-symbol formula give
\(\mathcal C_{\mathfrak p}(x,y)=(x,y)_{6,\mathfrak p}^{2}\) for all
local \(x,y\).  For rational \(a,b\), the product formulas for the chosen
torus lifts and for the Hilbert symbol therefore imply
\[
 \mathcal C_S(a,b)
 =\prod_{\mathfrak p\notin S}\mathcal C_{\mathfrak p}(a,b)^{-1}
 =H_S(a,b)^2.
\]
Both sides are continuous on \(F_S^\times\times F_S^\times\), so weak
approximation gives
\begin{equation}\label{eq:reflection-S-cocycle}
 \mathcal C_S(x,y)=H_S(x,y)^2
 \qquad(x,y\in F_S^\times).
\end{equation}
In particular this factor depends only on the classes in the fixed finite
group \(\mathcal V=F_S^\times/F_S^{\times6}\).  The local sections here
are the same ones used by the rational splitting.

\subsection{Bounds and a prescribed finite Whittaker function}

For \(v\mid\infty\), define the derivative in the logarithm of the norm by
\[
 \mathcal D_v f(y)
  =\left.\frac{\dd}{\dd t}f(e^{t/2}y)\right|_{t=0}.
\]
The factor \(1/2\) makes
\(|e^{t/2}y|_v=e^t|y|_v\).

We give both the Bessel calculation and bounds for the fixed rotations
of the vectors that we will use.  Choose complex coordinates in which
\(\psi_v(z)=e^{i c_v\operatorname{Re}z}\), with \(c_v>0\); a harmless
rotation and a change of sign give this form.  Put \(\varrho=|y|_v^{1/2}\).
In the dominant normalized principal series set
\(\nu_0=1/2+1/6=2/3\).  Its spherical section, restricted to
\(wn(x)a(y)\), is a constant times
\[
 \varrho^{2\nu_0}(\varrho^2+|x|_{\mathrm{ord}}^2)^{-2\nu_0},
\]
where \(|x|_{\mathrm{ord}}\) is the ordinary complex modulus.  The
Jacquet integral is absolutely convergent because \(2\nu_0>1\).
Using the Laplace representation of the negative power and the
two-dimensional Gaussian Fourier integral gives
\begin{align*}
 \varrho^{2\nu_0}\int_{\R^2}
   \frac{e^{-ic_v x_1}}{(\varrho^2+x_1^2+x_2^2)^{2\nu_0}}
     \,\dd x_1\dd x_2
 &=\frac{\pi \varrho^{2\nu_0}}{\Gamma(2\nu_0)}
     \int_0^\infty t^{2\nu_0-2}
       e^{-t\varrho^2-c_v^2/(4t)}\,\dd t\\
 &=c'_v \varrho K_{2\nu_0-1}(c_v\varrho)
  =c'_v \varrho K_{1/3}(c_v\varrho),\qquad c'_v\ne0.
\end{align*}
The Laplace and Gaussian integrals may be interchanged absolutely:
after integrating the Gaussian without its phase the integrand at
zero is \(O(t^{2\nu_0-2})\), with exponent greater than \(-1\).
This proves \eqref{eq:theta-complex-bessel}.

For an upper root translate, Whittaker equivariance gives
\begin{equation}\label{eq:theta-lie-multiplier}
 W_{\pi(n(z))\varphi}(a(y))=\psi_v(yz)W_\varphi(a(y)).
\end{equation}
Suitable complex linear combinations of the two real Lie derivatives
in \(z\) therefore multiply the spherical function by \(y\) and by
\(\bar y\), respectively.  Write
\(\omega_m(y)=(y/|y|_v^{1/2})^m\) for \(m\in\{1,-1\}\).
After multiplying each derivative vector by a nonzero constant, its
Whittaker function is
\begin{equation}\label{eq:theta-complex-derivative}
 W_{v,m}(a(y))
   =|y|_v\,\omega_m(y)K_{1/3}(c_v|y|_v^{1/2}).
\end{equation}

Here are uniform bounds for every vector in a fixed finite list
consisting of these derivative vectors and their fixed compact
rotations.  A rotation of a first Lie derivative of the spherical
vector is another first Lie derivative, since the spherical vector
is compact invariant and
\(\pi(k)\dd\pi(X)=\dd\pi(\operatorname{Ad}(k)X)\pi(k)\).
For a fixed real Lie algebra element \(X\), differentiate the
spherical section using
\[
 f^\circ(g)=
 \left(\frac{|\det g|_{\mathrm{ord}}}
 {|g_{21}|_{\mathrm{ord}}^2+|g_{22}|_{\mathrm{ord}}^2}\right)^{2\nu_0}.
\]
On \(wn(x)a(y)\), its first Lie derivative and any number of
\(\mathcal D_v\)-derivatives are the preceding spherical integrand
times a finite sum of bounded homogeneous rational functions of
\((y,x)\).  More precisely, for every norm derivative order \(b\)
and every \(x_1\)-derivative order \(h\), the resulting derivative is
bounded by
\[
 C_{b,h} \varrho^{2\nu_0}
 (\varrho^2+|x|_{\mathrm{ord}}^2)^{-2\nu_0-h/2}.
\]
This follows by differentiating the displayed quotient: each
derivative adds a polynomial numerator of the same degree as its
denominator, and an \(x_1\)-derivative lowers the homogeneous degree
by one.  The integral of the bound with \(h=0\) is
\(O(\varrho^{2-2\nu_0})=O(\varrho^{2/3})\).  For \(\varrho\ge1\), integrating by parts
\(h\) times against \(e^{-ic_vx_1}\) instead gives
\(O_{b,h}(\varrho^{2/3-h})\).  These integrable derivative bounds justify the
integrations by parts by cutoff and limit.  Consequently, for every
integer \(j\ge0\) and every \(A>0\), the fixed list satisfies
\begin{equation}\label{eq:theta-complex-bounds}
 |\mathcal D_v^j W(a(y))|
 \le C_{A,j}\min\bigl(|y|_v^{1/3},|y|_v^{-A}\bigr).
\end{equation}

At a finite place the corresponding elementary bounds hold for every
fixed smooth vector.  To be explicit, let
\(W_\varphi(y)=\lambda_v(\pi_v(a(y))\varphi)\).  There is an additive
open subgroup \(J\subset F_v\) for which \(n(J)\) fixes \(\varphi\).
The identity
\[
 W_\varphi(y)=\psi_v(yx)W_\varphi(y)\qquad(x\in J)
\]
forces \(W_\varphi(y)=0\) when \(|y|_v\) is sufficiently large.
For the other end, put \(Q_v=\#(\mathcal O_v/\mathfrak p_v)\) and
\(z_v=a(\varpi_v^3)\).  This element is central in the covering
torus.  By the two-cell Jacquet calculation
\cite[Proposition~I.2.1, p.~61, and Theorem~I.2.9(e), p.~72]{KP84}, the two
constituents of the unnormalized ordinary Jacquet module of the
principal series have \(z_v\)-eigenvalues of absolute values
\[
 Q_v^{-1}\quad\hbox{and}\quad Q_v^{-2}.
\]
These are the exponents \(1/2-1/6\) and \(1/2+1/6\) on a triple
valuation step.  Exactness of the ordinary Jacquet functor gives the
same possible eigenvalues on the exceptional subquotient.  Hence a
polynomial all of whose roots have those absolute values annihilates
the action of \(z_v\) on this Jacquet module.

For completeness, this gives a recurrence on the Whittaker values
themselves.  If \(P\) is such an annihilating polynomial, then
\[
 P(\pi_v(z_v))\varphi
   =\sum_{i=1}^h\bigl(\pi_v(n(x_i))-1\bigr)\varphi_i
\]
for finitely many \(x_i,\varphi_i\).  Apply
\(\lambda_v\pi_v(a(y))\).  For all sufficiently small \(|y|_v\) the
right side is zero because every \(\psi_v(yx_i)\) is one.  Cube
cocycles are trivial, so the left side is the recurrence \(P(E)W_\varphi(y)\),
where \(E W_\varphi(y)=W_\varphi(y\varpi_v^3)\).
On each valuation progression its solutions are sums of a polynomial
in the valuation times the powers of the roots of \(P\).
The unit variable ranges over only finitely many classes after
restricting to a sufficiently small unit subgroup which fixes the
vector and consists of cubes.  It follows that, for every fixed
\(\delta<1/3\),
\begin{equation}\label{eq:theta-finite-bounds}
 W_\varphi(y)=0\quad(|y|_v>C_\varphi),\qquad
 |W_\varphi(y)|\le C_{\varphi,\delta}|y|_v^\delta
       \quad(|y|_v\le1).
\end{equation}
All constants in this paragraph may depend on the place, the vector,
and the functional.  For a fixed finite list one may choose a common
positive \(\delta\).  The local twisted Jacquet quotients of
\(\Theta_v\) are finite dimensional by the quotient argument in
\Cref{subsec:theta-global-normalization}.  Thus the same conclusions
hold term by term for the finite sum of products at \(S\) in
\Cref{prop:theta-input}.

\begin{lemma}[Finite torus cutoff]\label{lem:finite-whittaker-cutoff}
Let \(S_f\) be a finite set of nonarchimedean places and let
\(\lambda_f\ne0\) be a Whittaker functional with character
\(\prod_{v\in S_f}\psi_v\) on a smooth representation of the product
of the local covers.  For a vector \(\varphi\), put
\[
 W_\varphi(y)=\lambda_f\!\left(\pi\bigl((a(y_v))_{v\in S_f}\bigr)
                         \varphi\right),
 \qquad y\in\prod_{v\in S_f}F_v^\times.
\]
Given any open subgroup \(H_f\) of the product of the unit groups,
there are a compact open product subgroup \(U\subset H_f\) and a
smooth vector \(\varphi_U\) for which
\[
 W_{\varphi_U}(y)=\one_U(y)
 \qquad\left(y\in\prod_{v\in S_f}F_v^\times\right).
\]
\end{lemma}

\begin{proof}
Choose a vector with \(\lambda_f(\varphi)\ne0\).  The map
\(y\mapsto W_\varphi(y)\) is locally constant near one.  Choose a
product of sufficiently small groups \(1+\mathfrak p_v^{m_v}\),
contained in \(H_f\), on which it is the constant
\(\lambda_f(\varphi)\); call that product \(U\).  As a subset of the
additive product \(\prod F_v\), \(U\) is compact and open.
By additive Fourier inversion choose a Schwartz--Bruhat function
\(\phi\) with
\[
 \int_{\prod F_v}\phi(t)\prod_v\psi_v(y_vt_v)\,\dd t=\one_U(y).
\]
The vector \(\varphi'=\int\phi(t)\pi(n(t))\varphi\,\dd t\) is a
well-defined smooth vector: its integrand is locally constant on the
compact support of \(\phi\), so the integral is a finite linear
combination of vectors.  The split root relation gives
\[
 W_{\varphi'}(y)
 =W_\varphi(y)\int\phi(t)\prod_v\psi_v(y_vt_v)\,\dd t
 =W_\varphi(y)\one_U(y).
\]
Divide \(\varphi'\) by \(\lambda_f(\varphi)\).
\end{proof}

Here is the choice at \(S\) that will be used later.  Write
\(S_f=S\setminus S_\infty\).  Let \(\mathscr L\) be any fixed finite
family of smooth unitary characters \(\mathcal L_S\) of
\(F_S^\times\), each trivial on \(U_S\), radially trivial at infinity,
and with angular component \(\omega_{m_v}\), \(m_v\in\{1,-1\}\), at
each complex place.  This permits the arithmetic character \(L_S\)
and its finitely many finite-order twists used below.  Factor the
nonzero global functional at its complex places, leaving a nonzero
functional at the finite places of \(S\).  Apply
\Cref{lem:finite-whittaker-cutoff} with an open subgroup contained in
the sixth powers and in the intersection of the kernels of the finite
components of all \(\mathcal L_S\in\mathscr L\).  We obtain one product
\(U=\prod_{v\in S_f}U_v\) and one finite vector.  For each
\(\mathcal L_S\), choose at each complex place the derivative of its
angular type.  The functions in \eqref{eq:theta-complex-derivative}
are nonzero at one, and their normalization gives an \(S\)-Whittaker
function
\begin{equation}\label{eq:theta-chosen-S-function}
 W_{S,\mathcal L_S}(a(y))=\one_U(y_f)
   \prod_{v\mid\infty}|y_v|_v\,\omega_{m_v}(y_v)
                      K_{1/3}(c_v|y_v|_v^{1/2}).
\end{equation}
In particular \(\mathcal L_S(y)^{-1}W_{S,\mathcal L_S}(a(y))\) is the
product of the same radial Bessel factors and \(\one_U(y_f)\) for
every member of the family.  The radial profile and its Gaussian
averaging kernel below may therefore be chosen in common for this
finite family.

The complex constant term is radial by \Cref{prop:theta-input}, whereas
\(\mathcal L_S^{-1}\) has a nontrivial angular component.  We will use this
distinction in \Cref{sec:reflection}, after the actual weighted theta sums
have been shown to descend to the quotient on which they are integrated.

\subsection{Integration on norm sheets}

Recall that \(r=\#S_\infty=[F:\Q]/2\), and define
\[
 X_u=\{y\in F_S^\times:|y|_S=u\}\qquad(u>0).
\]
For \(d\in F_S^\times\), put \(q_d=|d|_S\).  Multiplication by \(d\)
sends \(X_l\) to \(X_{q_dl}\), and inversion sends \(X_l\) to \(X_{1/l}\).
The forward and reflected theta terms will therefore lead to sheet
integrals at norms \(q_dl\) and \(q_d/l\), respectively.  We now bound
such one-variable profiles, including absolute envelopes that will justify
the ideal sums and the radial average.  The two signed orbit substitutions
will be proved with their actual integrands in \Cref{sec:reflection}.

There is at least one complex place.  Use the norm section \(\ell_u\) from
\Cref{sec:arithmetic}; it satisfies \(|\ell_u|_S=u\) and
\(\ell_u\ell_v=\ell_{uv}\).  Choose multiplicative Haar measures and Haar
measure on \(X_1\) so that
\begin{equation}\label{eq:sheet-disintegration}
 \int_{F_S^\times}f(y)\,\dd^\times y
 =\int_0^\infty\int_{X_u}f(y)\,\dd^\times_u y\,\frac{\dd u}{u},
\end{equation}
where the measure on \(X_u\) is the translate by \(\ell_u\) of that
on \(X_1\).  At a complex place we take angle measure of total mass
one, so its multiplicative measure is a fixed constant times
\(\dd t/t\) times angle measure, with \(t=|y_v|_v\).
At a finite place every valuation shell has the same multiplicative
measure.

\begin{lemma}[Norm-sheet profiles]\label{lem:sheet-profile}
Let \(\mathcal I\) be a finite index set.  For every
\(\kappa\in\mathcal I\), let \(f_\kappa\) be a finite sum of products
\[
 f_\kappa(y)=\sum_{\ell=1}^{L_\kappa}
                \prod_{v\in S}f_{\kappa,\ell,v}(y_v),
 \qquad y\in F_S^\times.
\]
Assume that for one \(\delta>0\) all these local factors satisfy the
following conditions.

\begin{enumerate}[label=\textup{(\alph*)}]
\item At a finite \(v\), \(f_{\kappa,\ell,v}\) is locally constant,
  vanishes for \(|y_v|_v>C_v\), and is
  \(O(|y_v|_v^\delta)\) for \(|y_v|_v\le1\).
\item At a complex \(v\), the factor is smooth on \(F_v^\times\) and, for
  every \(j\ge0\) and \(A>0\),
  \[
   |\mathcal D_v^j f_{\kappa,\ell,v}(y_v)|
      \le C_{A,j}\min(|y_v|_v^\delta,|y_v|_v^{-A}).
  \]
\end{enumerate}
All constants may depend on the finite family.  Define
\[
 H_\kappa(u)=\int_{X_u}f_\kappa(y)\,\dd^\times_u y,\qquad
 D=u\frac{\dd}{\dd u}.
\]
Fix one complex place \(v_0\), and define the nonnegative absolute
derivative profile
\[
 H^{\mathrm{abs}}_{\kappa,j}(u)
 =\sum_{\ell=1}^{L_\kappa}\int_{X_u}
   |\mathcal D_{v_0}^j f_{\kappa,\ell,v_0}(y_{v_0})|
   \prod_{v\ne v_0}|f_{\kappa,\ell,v}(y_v)|\,\dd^\times_u y.
\]
Then \(H_\kappa\) is smooth, and for every \(j\ge0\) and
\(\sigma>-\delta\),
\begin{equation}\label{eq:sheet-weighted-bound}
 |D^jH_\kappa(u)|\le H^{\mathrm{abs}}_{\kappa,j}(u)
      \le C_{\sigma,j}u^{-\sigma}\qquad(u>0).
\end{equation}
In particular, for any \(0<\alpha<\delta\) and every \(A>0\),
\begin{equation}\label{eq:sheet-two-ended-bound}
 |D^jH_\kappa(u)|+H^{\mathrm{abs}}_{\kappa,j}(u)
    \le C_{A,j}\min(u^\alpha,u^{-A}).
\end{equation}
The conclusions are unchanged if each \(f_\kappa\) is multiplied by
a fixed smooth unitary character of \(F_S^\times\) and by a bounded function
of its class in the finite group \(\mathcal V\).

Let \(\mathcal K:(0,\infty)\to\C\) have every absolute
moment
\[
 M_\beta(\mathcal K)=\int_0^\infty|\mathcal K(T)|T^\beta
                          \,\frac{\dd T}{T}<\infty
 \qquad(\beta\in\R).
\]
Both convolutions
\[
 \int_0^\infty\mathcal K(T)H_\kappa(uT)\,\frac{\dd T}{T},
 \qquad
 \Phi_\kappa(u):=\int_0^\infty\mathcal K(T)H_\kappa(u/T)
                                      \,\frac{\dd T}{T}
\]
and their norm derivatives obey \eqref{eq:sheet-two-ended-bound},
with constants depending also on the indicated moments.  For every
\(j\), the two absolute majorants
\(\int|\mathcal K(T)|H^{\mathrm{abs}}_{\kappa,j}(uT^{\pm1})\,\dd T/T\)
satisfy the same two-ended size bound.
For any \(x_i>0\), complex \(a_i\), indices
\(\kappa_i\in\mathcal I\), \(\tau\in\{1,-1\}\), and
\(\sigma>-\delta\), one has the absolute bound
\begin{equation}\label{eq:sheet-absolute-sum}
\begin{split}
 &\int_0^\infty|\mathcal K(T)|
       \sum_i|a_i|H^{\mathrm{abs}}_{\kappa_i,0}(u x_iT^\tau)
                                     \,\frac{\dd T}{T}\\
 &\hspace{12mm}\le
 C_{\sigma,0}u^{-\sigma}M_{-\tau\sigma}(\mathcal K)
       \sum_i|a_i|x_i^{-\sigma}.
\end{split}
\end{equation}
In particular the sum and integral may be interchanged if the final
series is finite.
\end{lemma}

\begin{proof}
For a finite place, let \(A_v(Q_v^{-k})\) be the supremum of the
absolute local factor on its \(k\)-th valuation shell.  Condition
(a) implies, for every \(\sigma>-\delta\),
\[
 \sum_{k\in\Z} A_v(Q_v^{-k})Q_v^{-\sigma k}<\infty:
\]
there are only finitely many negative \(k\), and the positive tail
is geometric with ratio \(Q_v^{-(\delta+\sigma)}\).
At a complex place condition (b) likewise gives
\[
 \int_{F_v^\times}|f_v(y)|\,|y|_v^\sigma\,\dd^\times y<\infty,
 \qquad
 \sup_{y\in F_v^\times}|y|_v^\sigma
                     |\mathcal D_v^j f_v(y)|<\infty.
\]
These statements remain valid for all factors in the fixed finite
family.

Use \(t_{v_0}=|y_{v_0}|_{v_0}\) as the dependent norm in
\eqref{eq:sheet-disintegration}.  On \(X_u\),
\[
 t_{v_0}=u\prod_{v\ne v_0}|y_v|_v^{-1},
\]
and the sheet measure is the product of the angle measure at \(v_0\)
and the multiplicative measures at the other places, up to its fixed
normalization.  Multiplying the absolute integrand by
\(u^\sigma=\prod_v|y_v|_v^\sigma\), take the weighted supremum of the
\(v_0\)-factor and the weighted integrals of all the others.  This
gives \eqref{eq:sheet-weighted-bound}, including its absolute
version.  On this parametrization \(D\) differentiates only the
\(v_0\)-factor and is exactly \(\mathcal D_{v_0}\).
The same weighted majorants for each derivative justify
differentiation under the integral on compact \(u\)-intervals.
Taking \(\sigma=-\alpha\) and then \(\sigma=A\) proves
\eqref{eq:sheet-two-ended-bound}.

At a complex place the sixth-power quotient is trivial, and at a
finite place it is finite.  Thus a bounded function of
\(\mathcal V\) is a finite linear combination of products of local
class indicators.  A unitary character factors over the finite
product \(S\); at a complex place its norm derivatives only introduce
fixed constants.  Multiplying by these functions preserves the
hypotheses after splitting into a finite family.

For the convolutions, apply the bound with \(\sigma=-\alpha\) or
\(\sigma=A\) under the \(T\)-integral.  For example the two bounds
for the reflected convolution have additional factors
\(M_{-\alpha}(\mathcal K)\) and \(M_A(\mathcal K)\).
The corresponding absolute derivative profiles dominate
differentiation on compact \(u\)-intervals.  Finally apply
\eqref{eq:sheet-weighted-bound} to each summand in
\eqref{eq:sheet-absolute-sum}; Tonelli and the definition of the
moment give the displayed bound.
\end{proof}

The lemma also gives the smooth separation on norm dyads that we
will use.  If \(\rho\) is a fixed smooth compactly supported
function on \(\R\), integration by parts twice in its ordinary
Fourier transform gives
\begin{equation}\label{eq:sheet-fourier-separation}
 \left\|\widehat{\rho(t)\Phi_\kappa(u e^t)}\right\|_{L^1(\R)}
 \le C_{A,\rho}\min(u^\alpha,u^{-A}).
\end{equation}
Indeed the \(L^1\) norm is bounded by a constant times the suprema
of the function and its derivatives of orders one and two;
these derivatives are controlled by
\eqref{eq:sheet-two-ended-bound}.  Fourier inversion with
\(t=\log x-\log y\), for \(x,y\) in fixed ratio intervals, therefore
separates the two norm variables with bounded \(L^1\) cost and
retains any chosen large-\(u\) power saving.

\subsection{Exact radial smoothing}

For the chosen \(S\)-function \eqref{eq:theta-chosen-S-function},
the angular characters cancel against \(\mathcal L_S^{-1}\).  Define its
radial profile by
\begin{equation}\label{eq:theta-V-profile}
 V(u)=\int_{X_u}\one_U(y_f)
       \prod_{v\mid\infty}|y_v|_v
           K_{1/3}(c_v|y_v|_v^{1/2})\,\dd^\times_u y.
\end{equation}
It is a positive, nonzero profile covered by
\Cref{lem:sheet-profile}.  In this subsection the Mellin convention is
\(\widehat f(s)=\int_0^\infty f(u)u^s\,\dd u/u\).
The unreflected theta average will contain \(V(q_{\mathfrak A}l)\).
At \(l=T/Z\), the next lemma averages this profile to exactly
\(P_G(q_{\mathfrak A}/Z)\).  After the Weyl operation the norm variable
is inverted, and the same kernel instead produces the controlled profiles
\(\int\mathcal K(T)H_\kappa(u/T)\,\dd T/T\).

\begin{lemma}[Exact Gaussian smoothing]\label{lem:gaussian-smoothing}
For \(V\) in \eqref{eq:theta-V-profile}, there are constants
\(C_0>0\) and \(A_0>0\) such that
\begin{equation}\label{eq:theta-V-mellin}
 \widehat V(s)
 =C_0 A_0^s
       \bigl[\Gamma(s+5/6)\Gamma(s+7/6)\bigr]^r
 \qquad(\re s>-5/6).
\end{equation}
There is a smooth function \(\mathcal K:(0,\infty)\to\C\) such that,
for every \(N>0\) and \(j\ge0\),
\begin{equation}\label{eq:theta-kernel-bounds}
 \left|(T\partial_T)^j\mathcal K(T)\right|
       \le C_{N,j}\min(T^N,T^{-N}),
\end{equation}
and, for every \(u>0\),
\begin{equation}\label{eq:theta-exact-gaussian}
 \int_0^\infty\mathcal K(T)V(uT)\,\frac{\dd T}{T}
  =P_G(u)=\frac{1}{2\sqrt{\pi}}
       \exp\!\left(-\frac{(\log u)^2}{4}\right).
\end{equation}
Its Mellin transform satisfies the exact sign convention
\begin{equation}\label{eq:theta-kernel-mellin}
 \widehat{\mathcal K}(-s)
       =\frac{e^{s^2}}{\widehat V(s)}.
\end{equation}
The quotient on the right is interpreted by its entire continuation.
For every fixed finite family of rotated \(S\)-profiles \(H_\kappa\)
covered by \Cref{lem:sheet-profile}, the functions
\(\Phi_\kappa(u)=\int\mathcal K(T)H_\kappa(u/T)\,\dd T/T\)
obey the signed and absolute bounds of that lemma.
\end{lemma}

\begin{proof}
The disintegration \eqref{eq:sheet-disintegration} and Tonelli give
the Mellin transform of \(V\) as the product of its local Mellin
integrals.  Each finite factor is the positive volume of \(U_v\),
because \(|y_v|_v=1\) there.  At a complex place put \(t=|y_v|_v\).
The classical Bessel integral, or the same Laplace integral used
above, gives
\begin{align}
 \int_0^\infty t^{s+1}K_{1/3}(c_v\sqrt t)\,\frac{\dd t}{t}
 &=2^{2s+1}c_v^{-2s-2}
       \Gamma(s+5/6)\Gamma(s+7/6),
       \qquad \re s>-5/6.                         \label{eq:theta-bessel-mellin}
\end{align}
The angle integrals only change the positive measure constant.
Multiplication over the \(r\) complex places proves
\eqref{eq:theta-V-mellin}.

Set
\[
 \mathcal R_V(s)=\frac{e^{s^2}}{\widehat V(s)}
   =C_0^{-1}A_0^{-s}e^{s^2}
       \bigl[\Gamma(s+5/6)\Gamma(s+7/6)\bigr]^{-r}.
\]
The reciprocal gamma function is entire, so \(\mathcal R_V\) is
entire.  On each fixed closed vertical strip, Stirling's formula
implies
\[
 |\mathcal R_V(\sigma+it)|
   \le C(1+|t|)^C e^{-t^2+\pi r|t|}.
\]
The same bound with a different polynomial holds after
multiplication by any power of \(s\).  Hence
\begin{equation}\label{eq:theta-kernel-definition}
 \mathcal K(T)=\frac{1}{2\pi i}\int_{(\sigma)}
                   \mathcal R_V(s)T^s\,\dd s
\end{equation}
converges absolutely and is independent of the real line
\(\sigma\): the horizontal sides of a rectangle tend to zero by the
Gaussian bound.  Differentiation in \(\log T\) inserts a power of
\(s\).  Moving the line to \(N\) for \(0<T\le1\), and to \(-N\)
for \(T\ge1\), proves \eqref{eq:theta-kernel-bounds}.
In particular every \(M_\beta(\mathcal K)\) in
\Cref{lem:sheet-profile} is finite.

Fourier inversion in \(\log T\) on a vertical line in
\eqref{eq:theta-kernel-definition} gives
\(\widehat{\mathcal K}(-s)=\mathcal R_V(s)\) on that line.
Both sides are entire, so \eqref{eq:theta-kernel-mellin} holds
everywhere.  For example on any real line \(\sigma>0\), absolute
Mellin Fubini is valid because
\[
 \int_0^\infty|\mathcal K(T)|T^{-\sigma}\frac{\dd T}{T}<\infty,
 \qquad
 \int_0^\infty V(u)u^\sigma\frac{\dd u}{u}<\infty.
\]
It follows that the Mellin transform of the left side of
\eqref{eq:theta-exact-gaussian} is
\(\widehat V(s)\widehat{\mathcal K}(-s)=e^{s^2}\).
The function on the right of \eqref{eq:theta-exact-gaussian} has
the same transform, by the elementary Gaussian integral.
Mellin inversion proves the identity.  The final assertion follows
from the convolution part of \Cref{lem:sheet-profile}.
\end{proof}
